\documentclass[../../main.tex]{subfiles}
\begin{document}

{\bf [解]}
 
簡単のため$ c = \cos x, \, s = \sin x$，$I = \int_0^\pi x s e^x dx$ とおく．
$I$を部分積分によって変形すると
\begin{align}
I
 &= [x s e^x]_0^\pi - \int_0^\pi (s + xc) e^x dx = -\int_0^\pi s e^x dx - \int_0^\pi c x e^x dx \label{eq:1}
\end{align}
となる．第二項をさらに変形すると
\begin{align}
\int_0^\pi c x e^x dx 
 &= [c x e^x]_0^\pi - \int_0^\pi (c - s x) e^x dx \\
 &= -\pi e^\pi - \int_0^\pi c e^x dx + I \label{eq:2}
\end{align}
であるから\cref{eq:1}に代入して
\begin{align}
 I &= -\int_0^\pi s e^x dx + \pi e^\pi + \int_0^\pi c e^x dx - I \\
\therefore
 I &= \frac{1}{2}\left(\pi e^\pi + \int_0^\pi c e^x dx -\int_0^\pi s e^x dx\right) \label{eq:3}
\end{align}
となる．ここで
\begin{align}
 \begin{cases}
  A = \int_0^\pi s e^x dx \\
  B = \int_0^\pi c e^x dx
 \end{cases}
\end{align}
とおいて\cref{eq:3}に代入すると
\begin{align}
 I &= \frac{\pi e^\pi + B - A}{2} \label{eq:4}
\end{align}
となる．
$A,B$を具体的に計算すると
\begin{align}
 \begin{cases}
A = \dfrac{1}{2} [e^x(s-c)]_0^\pi = \dfrac{1}{2}(e^\pi + 1) \\[5pt]
B = \dfrac{1}{2} [e^x(c+s)]_0^\pi = \dfrac{1}{2}(-e^\pi - 1)  
 \end{cases}
\end{align}
だから\cref{eq:4}に代入して
\begin{align} 
 I = \frac{\pi e^\pi - e^\pi - 1}{2} = \frac{(\pi-1)e^\pi - 1}{2}  
\end{align}
である．

{\bf [解2]}
$I$を部分積分すると
\begin{align}
I &= [-c x e^x]_0^\pi + \int_0^\pi c e^x(x+1) dx \\
  &= \pi e^\pi + B + \int_0^\pi c e^x x dx  \label{eq:5}
\end{align}
である．第三項をさらに部分積分して変形すると
\begin{align}
\int_0^\pi c e^x x dx &= [s e^x x]_0^\pi - \int_0^\pi s e^x (x+1) dx = -A - I
\end{align}
だから\cref{eq:5}に代入すると
\begin{align}
  I &= \pi e^\pi + B - A - I \\[5pt]
\therefore 
  I &= \frac{\pi e^\pi + B - A}{2}
\end{align}
となって\cref{eq:4}に帰着する．（以下略）

\end{document}
